\section{Simulation Results}
\label{sec:results}
All 170 measured executions complete the frozen workload and pass the
independent raw-evidence audit. This section separates protection traffic,
recovery preparation and whole-workflow completion. Figures show all five
paired admission orders and 95\% intervals where applicable.
%RESULT_TABLE%

\subsection{Healthy protection cost}
%RESULT_HEALTHY%
\begin{figure*}[t]
\centering\includegraphics[width=\textwidth]{figures/aligned-healthy.pdf}
\caption{Healthy execution. (a) Full minus control mean WCT, aggregated within
each seed before comparison. (b) Bytes actually transferred for protection
across the complete workload, including work on later-aborted copies.
Traffic is in TiB ($2^{40}$ bytes). Open circle, square, up/down triangle
and plus markers identify seeds 0--4; filled diamonds and bars show means
and 95\% Student-$t$ intervals. Modeled zero latency cost would not imply zero physical
controller or kernel overhead.}
\Description{Paired healthy workflow-latency differences and protection
traffic for unprotected, DRAM-only and full treatments.}
\label{fig:healthy}
\end{figure*}

\subsection{Failure impact on complete workflows}
%RESULT_FAILURE%
\begin{figure*}[t]
\centering\includegraphics[width=\textwidth]{figures/aligned-failure.pdf}
\caption{Base-setting failure outcomes. (a) Descriptive WCT ECDFs over 160
workflows in each of five runs per arm; the 800 values are not independent
replicates. (b,c) Paired F1-minus-F0 penalty contrasts for the whole-workload
mean and the prespecified focus workflow. Negative values favor full.
Intervals use five admission orders, not workflow-level pseudo-replicates.}
\Description{Raw workflow completion distributions and controlled contrasts
against unprotected recomputation and DRAM-only protection.}
\label{fig:failure}
\end{figure*}

\subsection{Recovery state and preparation}
%RESULT_RECOVERY%
\begin{figure*}[t]
\centering\includegraphics[width=\textwidth]{figures/aligned-recovery.pdf}
\caption{The five prespecified focus requests. (a) DRAM/full bars partition each
prompt into HBM copy, DRAM load and recomputation. (b) Paired prefix-ready
times include queueing and preparation; lines connect treatments within a
seed. (c) Prefix-ready saving against the F0-corrected focus-WCT penalty
contrast; numbers identify seeds (1--3 nearly coincide). Seed 4 prepares
faster yet finishes its workflow later. No regression is fitted.}
\Description{Per-seed recovery-source proportions, paired prefix preparation times,
and their relationship to corrected workflow completion penalties.}
\label{fig:recovery}
\end{figure*}

\subsection{Mechanism ablations}
%RESULT_ABLATIONS%
\begin{figure*}[t]
\centering\includegraphics[width=\textwidth]{figures/aligned-ablations.pdf}
\caption{Ablation minus full: (a,b) failure penalties corrected by each arm's
own F0 run; (c) healthy copy traffic. Positive values favor full. Slack
sets the signal to zero; AIMD fixes the initial quota; LRU replaces retention.
All keep request ownership, copying and cleanup. Each estimate uses five
paired admission orders; the byte effect does not establish a WCT gain.}
\Description{Whole-workload, focus-workflow and healthy-byte effects of removing
Slack, fixing quota control, and replacing retention with LRU.}
\label{fig:ablations}
\end{figure*}

\subsection{Sensitivity to modeled conditions}
%RESULT_SENSITIVITY%
\begin{figure*}[t]
\centering\includegraphics[width=\textwidth]{figures/aligned-sensitivity.pdf}
\caption{Full minus DRAM failure-penalty contrasts across the declared
settings. $C$ is workflow concurrency and $\beta$ the serialized fraction of
modeled batch service. Two faults use their two prespecified focus workflows
and reuse base F0 controls. Open markers show all seeds; intervals are exploratory,
unadjusted comparisons. Negative values favor full.}
\Description{Paired sensitivity to concurrency, batch service composition,
path bandwidth and two sequential failures.}
\label{fig:sensitivity}
\end{figure*}
